\subsection{\texorpdfstring{The Grothendieck Group \(K_{0}(A)\)}{The Grothendieck Group K\_\{0\}(A)}}

Let \(\mathrm{A}\) be a ring. The set \(\mathcal{P}(\mathrm{A})\) of classes of finitely generated projective A-modules is additive; we denote the Grothendieck group \(\mathrm{K}(\mathcal{P}(\mathrm{A}))\) by \(\mathrm{K}_0(\mathrm{A})\) .

For the law of composition \((\mathrm{E},\mathrm{E}^{\prime})\mapsto\mathrm{cl}(\mathrm{E}\oplus\mathrm{E}^{\prime})\) , the set \(\mathcal{P}(\mathrm{A})\) is a commutative monoid. Moreover, every exact sequence of projective A-modules

\[
0 \longrightarrow \mathrm{E} ^ {\prime} \longrightarrow \mathrm{E} \longrightarrow \mathrm{E} ^ {\prime \prime} \longrightarrow 0
\]

splits (II, §2, No.~2, p.~231, Proposition 4), so that E is isomorphic to \(E^{\prime} \oplus E^{\prime\prime}\) . The mapping \(E \mapsto [E]\) from \(\mathcal{P}(A)\) to \(K_{0}(A)\) therefore defines an isomorphism from the group of differences of the monoid \(\mathcal{P}(A)\) to \(K_{0}(A)\) (VIII, p.~188).

For every finitely generated projective module P, there exists a finitely generated projective module \(P'\) such that \(P \oplus P'\) is free (II, §2, No.~2, p.~232, Corollary 1). Let E and F be finitely generated projective A-modules; by I, §2, No.~4, p.~18, Proposition 6, we have \([E] = [F]\) in \(K_{0}(A)\) if and only if there exists a finitely generated free A-module L such that the A-modules \(E \oplus L\) and \(F \oplus L\) are isomorphic. We then say that E and F are stably isomorphic;

this does not necessarily imply that E and F are isomorphic (VIII, p.~207, Exercise 2 and VIII, p.~210, Exercise 14).

When the ring A is commutative, there exists a commutative ring structure on the additive group \(K_{0}(A)\) whose multiplication is characterized by the formula \([\mathrm{E}]_{\mathcal{P}(\mathrm{A})}[\mathrm{F}]_{\mathcal{P}(\mathrm{A})} = [\mathrm{E} \otimes_{\mathrm{A}} \mathrm{F}]_{\mathcal{P}(\mathrm{A})}\) (VIII, p.~199, Proposition 10).

Remark. --- Let A be a semisimple ring. Every A-module is then semisimple and projective (VIII, p.~138, Proposition 4); we therefore have the equality

\[
\mathscr {L} \mathscr {F} (\mathrm{A}) = \mathscr {S S} (\mathrm{A}) = \mathscr {P} (\mathrm{A})
\]

(cf.~No.~4 for the definitions of \(\mathcal{LF}(\mathrm{A})\) and \(\mathcal{SS}(\mathrm{A})\) ). We therefore have \(\mathrm{K}_0(\mathrm{A}) = \mathrm{R}(\mathrm{A})\) by the definitions of these Grothendieck groups.

Example. --- If every finitely generated projective A-module is free, then the rank defines an isomorphism from \(K_{0}(A)\) to Z. This is, in particular, the case when A is a principal ideal domain (VII, §1, No.~3, p.~15, Corollary 3) or when A is a local ring (VIII, p.~36, Corollary 6).

